%! Polar Function Graphs — Unit 7, Lesson 7.7 — Exit Ticket (Answer Key)
\documentclass[10pt]{article}
\usepackage{precalculus-article}
\usepackage{precalculus-key}   % pulls in -boxes; do NOT also load -boxes
\newcommand{\TallMath}[1]{$\displaystyle #1\rule[-1.4em]{0pt}{3.2em}$}

\begin{document}

\pageheader{Unit 7, Lesson 7.7}{Exit Ticket --- Answer Key}
\namedateperiod

\vspace{0.12in}

\begin{notesbox}{Exit Ticket --- Key}

\textbf{1.}\quad Complete the table for $r = 3\cos\theta$.

\vspace{4pt}
\renewcommand{\arraystretch}{1.7}
\begin{tabular}{|c|c|c|c|c|}
\hline
$\theta$ & $0$ & $\dfrac{\pi}{2}$ & $\pi$ & $\dfrac{3\pi}{2}$ \\
\hline
$r = 3\cos\theta$ & \ans{3} & \ans{0} & \ans{$-3$} & \ans{0} \\[4pt]
\hline
\end{tabular}

\vspace{0.18in}
\textbf{2.}\quad $\theta = $ \ans{$\pi/2$ and $3\pi/2$}

\vspace{6pt}
\ans{The polar curve passes through the pole (origin) at those angles.}

\vspace{0.16in}
\textbf{3.}\quad

\begin{enumerate}[label=(\Alph*), leftmargin=2em, itemsep=6pt]
  \item $\theta$ is the output (radius) and $r$ is the input (angle).
  \item \textbf{$\theta$ is the input (angle) and $r$ is the output (radius).}
  \item Both $\theta$ and $r$ are inputs to the function.
  \item Both $\theta$ and $r$ are outputs of the function.
\end{enumerate}

\vspace{4pt}
Answer: \ans{(B)}

\begin{teachernote}
Common error: students choose (A) because they see $r$ written first in the notation
$r = f(\theta)$ and confuse the variable on the left side with the input.
Reinforce: in $r = f(\theta)$, $\theta$ is inside the function (the input) and
$r$ is the result (the output).
\end{teachernote}

\end{notesbox}

\end{document}
